\documentclass[11pt]{article}
\usepackage[a4paper,margin=1in]{geometry}
\usepackage{amsmath,amssymb}
\usepackage{hyperref}
\title{Reference Frames and Conventions in PyTex}
\author{PyTex Contributors}
\date{\today}

\begin{document}
\maketitle

\section{Frame Domains}

PyTex distinguishes at least the following frame domains:

\begin{itemize}
  \item crystal frame,
  \item specimen frame,
  \item map frame,
  \item detector frame,
  \item laboratory frame,
  \item reciprocal frame.
\end{itemize}

\section{Transform Semantics}

A frame transform is represented as a proper rigid transform
\[
  \mathbf{x}_{\text{target}} = \mathbf{R}\mathbf{x}_{\text{source}} + \mathbf{t},
\]
with $\mathbf{R} \in SO(3)$ and explicit source and target frames.

\section{Why This Matters}

The same numerical vector can mean different things in different software environments. PyTex aims to force those meanings to remain inspectable.

\end{document}
